% TOP_PROBLEM: {"id": 4700010, "title": "Number of centers", "category_id": 11, "category": {"id": 11, "name": "dynamical_systems", "display_name": "Dynamical Systems", "description": "Problems about long-term behavior of deterministic systems, Hamiltonian dynamics, and periodic orbits.", "slug": "dynamical-systems", "order_index": 11, "created_at": "2026-07-31T15:26:25.671Z"}, "status": "partially_solved", "problem_number": "AMR-046-0010", "set_id": 13}
% FIRST_POSTED: 2026-10-02
% ENTRY_KIND: statement_only
% UnsolvedMath ID 4700010; source catalogue code AMR-046-0010
% !TeX program = lualatex
% Definition and sources only; this file is not an LLM solution attempt.
\documentclass[11pt,a4paper]{article}
\usepackage[margin=25mm]{geometry}
\usepackage{fontspec}
\setmainfont{lmroman10-regular.otf}[BoldFont=lmroman10-bold.otf,
 ItalicFont=lmroman10-italic.otf,BoldItalicFont=lmroman10-bolditalic.otf]
\usepackage{amsmath,amssymb,mathtools}
\usepackage{unicode-math}
\setmathfont{latinmodern-math.otf}
\AtBeginDocument{\let\setminus\smallsetminus}
\usepackage{xurl,hyperref}
\hypersetup{colorlinks=true,urlcolor=blue}
\setcounter{secnumdepth}{0}
\setlength{\parindent}{0pt}
\setlength{\parskip}{5pt}
\setlength{\emergencystretch}{3em}
\begin{document}
\section*{Number of centers}\label{4700010}
{\small UnsolvedMath ID 4700010 · AMR-046-0010 · Dynamical Systems · AMR Open Problem Lists}
\subsection{Definitions and mathematical statement}
Let $X=(P,Q)$ be a real planar polynomial vector field of degree at most $n$,
meaning $\max(\deg P,\deg Q)\le n$. An equilibrium $z_0$ satisfies
$P(z_0)=Q(z_0)=0$. It is a center if it is isolated and has a punctured
neighborhood filled by nonconstant periodic orbits surrounding it.
This definition includes degenerate centers: no nonvanishing determinant
or prescribed pair of eigenvalues is assumed for $DX(z_0)$.

Define
\[
 \mathcal C_n=\sup_{\deg X\le n}
       \#\{z_0\in\mathbb R^2:z_0\text{ is a center of }X\}.
\]
The problem is to determine $\mathcal C_n$ for every integer $n\ge4$.
A complete determination must give both an upper bound for all real
coefficient choices and examples attaining that bound, or clarify whether
the supremum is finite and attained. Distinct centers are counted by their
locations, not by the number of periodic orbits in their surrounding annuli.

Common polynomial factors in $P$ and $Q$ are allowed, but non-isolated
equilibrium curves are not counted as centers. Counting only nondegenerate
linear centers would change the problem. Nor is the task the classification
of centers at one preassigned point: several different center regions can
coexist in a single vector field. The total degree restriction, rather than
separate bounds on each individual exponent, is the one used in defining
the sequence $\mathcal C_n$.
\subsection{Short English statement}
Find the greatest number of distinct centers that one degree-$n$ planar polynomial vector field can have, for $n\ge4$.
\subsection{Sources}
\begin{itemize}
\item[S1] ULAM.AI and the UnsolvedMath Contributors, supplied problems.json, record 4700010 (AMR-046-0010), AMR Open Problem Lists. Catalogue identity and source transcription; CC BY 4.0.
\url{https://huggingface.co/datasets/ulamai/UnsolvedMath}.
\item[S2] Gasull - Some open problems in low dimensional dynamical systems (2020), Problem environment 10: Number of centers. Original source indexed by AMR-046-0010.
\url{https://arxiv.org/abs/2012.02524}.
\end{itemize}
\end{document}
